\documentclass{article}
\usepackage[UTF8]{ctex}
\usepackage[a4paper,margin=2.2cm]{geometry}
\usepackage{amsmath,amssymb,mathtools}
\usepackage{hyperref}
\usepackage{bookmark}
\usepackage{lmodern}

\hypersetup{
  pdftitle={概率论 HW1},
  colorlinks=true,
  linkcolor=blue,
  urlcolor=blue
}

% % % % % % % % % % % % % % % % % % % % % % % % % % % % % % % %
% Common Macros for LaTeX
% % % % % % % % % % % % % % % % % % % % % % % % % % % % % % % %

% Useful packages
\usepackage{amsfonts}
\usepackage{amsmath}
\usepackage{amssymb}
\usepackage{amsthm}
\usepackage{enumerate}
\usepackage{graphicx}
\usepackage{multicol}
\usepackage[ruled,linesnumbered]{algorithm2e}

% Math operators and common symbols
\newcommand{\dif}{\mathrm{d}}
\newcommand{\innerProd}[1]{\left\langle #1 \right\rangle}
\newcommand{\difFrac}[2]{\frac{\dif #1}{\dif #2}}
\newcommand{\pdfFrac}[2]{\frac{\partial #1}{\partial #2}}
\newcommand{\T}{\mathrm{T}}
\newcommand{\norm}[1]{\left\| #1 \right\|}
\newcommand{\abs}[1]{\left| #1 \right|}

% Vector and Matrix shortcuts
\newcommand{\vecTwo}[2]{
  \begin{bmatrix}
    #1 \\
    #2
\end{bmatrix}}
\newcommand{\vecThree}[3]{
  \begin{bmatrix}
    #1 \\
    #2 \\
    #3
\end{bmatrix}}
\newcommand{\vecTwoH}[2]{
  \begin{bmatrix}
    #1 & #2
  \end{bmatrix}
}
\newcommand{\vecThreeH}[3]{
  \begin{bmatrix}
    #1 & #2 & #3
  \end{bmatrix}
}

% Common fields
\newcommand{\R}{\mathbb{R}}
\newcommand{\C}{\mathbb{C}}
\newcommand{\Q}{\mathbb{Q}}
\newcommand{\Z}{\mathbb{Z}}
\newcommand{\N}{\mathbb{N}}

% Common Operators
\renewcommand{\Re}{\mathrm{Re}}
\renewcommand{\Im}{\mathrm{Im}}
\newcommand{\Arg}{\mathrm{Arg}}
\newcommand{\Log}{\mathrm{Log}}

\newcommand{\arccot}{\mathrm{arccot}}
\newcommand{\arcsec}{\mathrm{arcsec}}
\newcommand{\arccsc}{\mathrm{arccsc}}

% Theorem-like environments
\theoremstyle{definition}
\newtheorem{theorem}{Theorem}[section]
\newtheorem{lemma}[theorem]{Lemma}
\newtheorem{corollary}[theorem]{Corollary}
\newtheorem{definition}[theorem]{Definition}
\newtheorem{remark}[theorem]{Remark}
\newtheorem{exercise}{Exercise}[section]

% Support for individual Chinese characters using fontspec (for LuaLaTeX)
\usepackage{fontspec}
\newfontfamily\zhfont{SimSun} % Search system fonts by name
\newcommand{\zhstr}[1]{{\zhfont #1}}

% Solutions and proofs
\AtBeginDocument{
  \ifdefined\ctexset
  \newenvironment{solution}{
  \begin{proof}[解]}{
  \end{proof}}
  \else
  \newenvironment{solution}{
  \begin{proof}[Solution]}{
  \end{proof}}
  \fi
}

\usepackage{luacode}
% Utility to get environment variables (requires lualatex)
\newcommand{\getEnv}[2]{\directlua{tex.print(os.getenv("#1") or "#2")}}

\newcommand{\name}{\getEnv{NAME}{NAME}}
\newcommand{\uid}{\getEnv{UID}{UID}}
\newcommand{\combin}[2]{\binom{#1}{#2}}

\newenvironment{problem}{\par\medskip\noindent\textbf{题目.}\ }{\par\smallskip}

\title{概率论 HW1}
\author{\name \uid}
\date{2026-09-20}

\begin{document}
\maketitle

\section*{第 5 题\quad 有放回摸球的次序}

\begin{problem}
  从盛有号码为 $1,2,\dots,N$ 的球的箱子中有放回地摸了 $n$ 次球，依次记其号码，求这些号码按严格上升次序排列的概率.
\end{problem}

\begin{solution}
  记第 $i$ 次摸到的号码为 $a_i$，则样本空间是 $\{1,\dots,N\}^n$，共有 $N^n$ 个样本点。所求事件为
  \[
    A=\{a_1<a_2<\cdots<a_n\}.
  \]
  $n>N$ 时 $n$ 个号码不可能互不相同，$P(A)=0$。设 $n\le N$。

  $A$ 中每个样本点的 $n$ 个号码互不相同；反过来，任取 $\{1,\dots,N\}$ 中的 $n$ 个号码，按从小到大排列就得到 $A$ 中一个样本点。所以 $A$ 与 $n$ 元子集一一对应，
  \[
    |A|=\binom Nn,
    \qquad
    P(A)=\frac{\binom Nn}{N^n}
    =\frac{N(N-1)\cdots(N-n+1)}{n!\,N^n}.
  \]

\end{solution}

\section*{第 9 题\quad 取鞋问题}

\begin{problem}
  从 $n$ 双不同的鞋子中任取 $2r$ ($2r\le n$) 只，求下列事件的概率:
  (1) 没有成双的鞋子; (2) 只有一双鞋子; (3) 恰有两双鞋子; (4) 有 $r$ 双鞋子.
\end{problem}

\begin{solution}
  总取法为 $\binom{2n}{2r}$。

  \begin{enumerate}[(1)]
    \item 没有成双：取到的鞋来自 $2r$ 双不同的鞋，每双取一只，共 $\binom{n}{2r}2^{2r}$ 种，
      \[
        P_1=\frac{\binom{n}{2r}2^{2r}}{\binom{2n}{2r}} .
      \]

    \item 只有一双：先定下整双取出的那一双，其余 $2r-2$ 只从另外 $n-1$ 双中取，每双取一只，共
      $\binom n1\binom{n-1}{2r-2}2^{2r-2}$ 种，
      \[
        P_2=\frac{\binom n1\binom{n-1}{2r-2}2^{2r-2}}{\binom{2n}{2r}} .
      \]

    \item 恰有两双：同理，
      \[
        P_3=\frac{\binom n2\binom{n-2}{2r-4}2^{2r-4}}{\binom{2n}{2r}} .
      \]

    \item 有 $r$ 双：$2r$ 只恰好配成 $r$ 双，就是取出的鞋来自某 $r$ 双且每双两只都取到，只需选出这 $r$ 双，
      \[
        P_4=\frac{\binom nr}{\binom{2n}{2r}} .
      \]
  \end{enumerate}

  这里约定 $\binom nk=0$（$k>n$ 或 $k<0$），$r=1$ 时 (3) 式即为 $0$。
\end{solution}

\section*{第 10 题\quad 电梯问题}

\begin{problem}
  10 层楼的一架电梯在底层登上 7 位乘客，从第二层起乘客可离开电梯，求每层至多一位乘客离开的概率（乘客在各层离开是等可能的）.
\end{problem}

\begin{solution}
  从第二层到第十层共 $9$ 层，7 位乘客各选一层离开，共有 $9^7$ 种等可能的选法。

  每层至多一人离开，就是 7 人选的楼层互不相同：先选出 $7$ 层，再把 7 人分派到这 $7$ 层，共有
  \[
    9\cdot8\cdot7\cdot6\cdot5\cdot4\cdot3=181440
  \]
  种，故
  \[
    P=\frac{9\cdot8\cdot7\cdot6\cdot5\cdot4\cdot3}{9^7}
    =\frac{2240}{59049}\approx 0.0379 .
  \]
\end{solution}

\section*{第 15 题\quad 船靠码头（会面问题）}

\begin{problem}
  某码头只能容纳一只船. 现知某日将独立来到两只船，且在 $24$ h 内来到的可能性均相等. 如果它们停靠的时间分别为 $3$ h 和
  $4$ h，求有一船要在江中等待的概率.
\end{problem}

\begin{solution}
  记停靠 $3$ h 的船在时刻 $X$ 到达，停靠 $4$ h 的船在时刻 $Y$ 到达，则 $(X,Y)$ 在正方形 $[0,24]^2$ 上均匀分布，面积为
  $576$。两只船都要靠码头，一只要等待当且仅当停留时段 $[X,X+3]$ 与 $[Y,Y+4]$ 相交，即
  \[
    X<Y+4,\qquad Y<X+3 .
  \]

  取对立事件。两船不相遇只有两种情形：$X\ge Y+3$，即 $3$ h 的船到得晚，$4$ h 的船已经走了；或 $Y\ge X+4$，即 $4$ h
  的船到得晚，$3$ h 的船已经走了。它们分别是正方形右下角（由直线 $x-y=3$ 截出，直角边长 $21$）和左上角（由
  $y-x=4$ 截出，直角边长 $20$）的三角形，面积为
  \[
    \frac{21^2}{2}=220.5,\qquad \frac{20^2}{2}=200,
  \]
  两区域不相交，故
  \[
    P\left(\text{有一船等待}\right)
    =\frac{576-220.5-200}{576}
    =\frac{155.5}{576}=\frac{311}{1152}\approx 0.270 .
  \]

\end{solution}

\section*{第 19 题\quad 投硬币与格线}

\begin{problem}
  在一张打方格子的纸上投一枚直径为 $1$ 的硬币，问方格要多小才能使硬币与线不相交的概率小于 $1\%$？
\end{problem}

\begin{solution}
  设方格边长为 $a$，硬币半径为 $\dfrac12$。固定一个方格来看，圆心在其中均匀分布，面积为 $a^2$。硬币不与格线相交，当且仅当圆心到方格四条边的距离都大于
  $\dfrac12$，即圆心落在方格中央边长为 $a-1$ 的小正方形内，所以
  \[
    p(a)=\frac{(a-1)^2}{a^2},\qquad a>1 .
  \]

  要求 $p(a)<0.01$，即
  \[
    \frac{a-1}{a}<0.1
    \iff 1-\frac1a<0.1
    \iff a<\frac{10}{9}\approx 1.111 .
  \]
  故所求非平凡范围为
  \[
    1<a<\frac{10}{9} .
  \]
  $a\le1$ 时硬币必与格线相交，不相交的概率为 $0$，属平凡情形。

\end{solution}

\end{document}
